\documentclass[11pt]{article}
\usepackage[no-math]{fontspec}
\usepackage[french]{babel}
\usepackage{amsmath}
\usepackage{amssymb}
\usepackage{amsthm}
\usepackage{mathrsfs}
\usepackage{tikz-cd}
\usepackage{enumitem}
\usetikzlibrary{babel}

\newcommand{\printedpage}[1]{\marginpar{\footnotesize p.~#1}}
\newcommand{\msp}[1]{\marginpar{\footnotesize [#1]}}
\newcommand{\HOT}{\mathrm{HOT}}
\newcommand{\Cat}{\mathrm{Cat}}
\newcommand{\Ob}{\mathrm{Ob}}
\newcommand{\Fl}{\mathrm{Fl}}
\newcommand{\id}{\mathrm{id}}
\newcommand{\cA}{\mathcal{A}}
\newcommand{\ul}[1]{\underline{\mathrm{#1}}}
\newcommand{\circled}[1]{\tikz[baseline=(c.base)]\node[draw,circle,inner sep=1pt](c){#1};}

\begin{document}

\printedpage{39}
% head: A. Grothendieck : Dérivateurs, Ch. VI, Édition des Universités de Montpellier II et Paris VII

\[
\begin{tikzcd}[column sep=large, row sep=large]
 & \HOT_W(S)\approx(\Cat/S)(W_S^{\mathrm{d}})^{-1} \\
\HOT^{\mathrm{lc}}_W(S)\approx(\Cat/S)W_S^{-1} \arrow[ur, "i"] \arrow[r, "C"] \arrow[dr, "j"'] & (\ul{Parf}\,S)W_S^{-1} \arrow[u, "\text{\footnotesize car $W_S=W_S^{\mathrm{d}}$ sur $\ul{Parf}$}"'] \arrow[d, "\text{\footnotesize car $W_S=W_S^{\mathrm{g}}$ sur $\ul{Parf}$}"] \\
 & \HOT^{o}_W(S)\approx(\Cat/S)(W_S^{\mathrm{g}})^{-1}.
\end{tikzcd}
\]

\msp{page 64}
On voit par là qu'on a des isomorphismes canoniques
\[
\rho i\simeq\id,\qquad\sigma j\simeq\id
\]
dans la catégorie des endofoncteurs de $(\Cat/S)W_S^{-1}$ ($\approx\HOT^{\mathrm{lc}}_W(S)$), isomorphisme qui s'explicite ici par le fait que $\rho i$, $\sigma j$ sont déduits du même foncteur $C$ par passage aux localisés $W_S^{-1}$ (donc en fait $\rho i=\sigma j$ !), et qu'on a une flèche fonctorielle en $X$
\[
X\xrightarrow{i_X}C(X),
\]
qui est dans $W$ pour tout $X$. Ainsi, $\rho$ et $\sigma$ sont à isomorphisme près des \textit{rétractions} des catégories $\HOT_W(S)$, $\HOT^{o}_W(S)$ sur la catégorie $\HOT^{\mathrm{lc}}_W(S)$. De plus, ces rétractions sont ce qu'on appelle des foncteurs de localisation (comme on voit sur le diagramme de la page 62). (On dit que $C\xrightarrow{F}C'$ est un \textit{foncteur de localisation} si, désignant par $\Sigma\subseteq\Fl\,C$ l'ensemble des flèches de $C$ transformées en des isomorphismes par $F$, le foncteur
\[
C\Sigma^{-1}\xrightarrow{\bar{F}}C'
\]
qui factorise $F$ est une équivalence de catégories.)

\msp{page 65}
On sait d'ailleurs que si un foncteur $f$ a un adjoint $g$ (à gauche ou à droite), alors $f$ est un foncteur \texttt{[de]} localisation si et seulement si l'adjoint $g$ est pleinement fidèle. Pour prouver que $i$, $j$ sont pleinement fidèles, il serait donc tentant de prouver qu'ils sont adjoints des foncteurs $\rho$, $\sigma$.

Pour nous y reconnaître, je vais poser
\[
\left\{
\begin{aligned}
&\cA=\Cat/S,\ \ \Sigma_0=W_S^{\mathrm{d}},\ \ \Sigma=W_S\qquad\quad(\subseteq\Fl(\cA))\\
&C:\cA\longrightarrow\cA\ \text{ le foncteur familier, qui a la vertu}\\
&\qquad\qquad\qquad\quad C(\Sigma)\subseteq\Sigma_0\quad.
\end{aligned}
\right.
\]

On a donc la situation
\printedpage{40}
% head: A. Grothendieck : Dérivateurs, Ch. VI, Édition des Universités de Montpellier II et Paris VII
\[
(*)\qquad\qquad(\cA,\Sigma_0)\ \underset{C}{\overset{\rho_0\,=\,\id_{\cA}}{\rightleftarrows}}\ (\cA,\Sigma),
\]
induisant par passage aux localisés
\[
\cA\Sigma_0^{-1}\quad\underset{\bar{C}}{\overset{\rho\,=\,\bar{\rho}_0}{\rightleftarrows}}\ \cA\Sigma^{-1},
\]
on veut vérifier que $\rho$ et $C$ \texttt{[plutôt $\bar{C}$]} sont adjoints l'un de l'autre, $C$ \texttt{[plutôt $\bar{C}$]} \textit{sûrement à droite,} ce qui signifie que l'on doit avoir des homomorphismes fonctoriels
\[
\left\{
\begin{aligned}
&\id_{\cA\Sigma^{-1}}&&\rightleftarrows&&\rho\,\bar{C}\quad(=\overline{\rho_0C})\\
&\id_{\cA\Sigma_0^{-1}}&&\rightleftarrows&&\bar{C}\rho\quad(=\overline{C\rho_0})
\end{aligned}
\right.
\]
(\footnote[41]{mais si $C$ est à \textit{droite,} les flèches ici doivent aller en sens inverse et c'est bien à quoi je dois me résoudre p.~67.}) \texttt{[les flèches en bas vont en direction corrigée, mais les flèches en haut ne sont pas effacées]} satisfaisant deux relations de compatibilité sur lesquelles on reviendra. On peut espérer trouver ces homomorphismes fonctoriels à partir de
\[
\left\{
\begin{aligned}
&\id_{\cA}&&\xrightarrow{\ i\ }&&\rho_0C&&=&&C\\
&\id_{\cA}&&\xleftarrow{\ p\ }&&C\rho_0&&=&&C
\end{aligned}
\right.
\]

\msp{page 66}
et dans notre cas on dispose bel et bien de
\[
i:\id_{\cA}\longrightarrow C\qquad\text{i.e. de }\ i_X:X\longrightarrow C(X)
\]
fonctoriel en $X\in\Ob\,\cA$ (et de plus $i_X$ est dans $\Sigma$). Il n'y a pas cependant de flèche naturelle en sens inverse $p_X$
\[
C(X)\xrightarrow{\ ?\ }X.
\]
J'avais songé à la projection $\pi_X:C(X)\longrightarrow X$, déduite de la construction $C(X)\simeq X\times_S(\ul{Ch}_\infty(S),\sigma)$, mais ce n'est malheureusement pas un $S$-morphisme. Comme on cherche une flèche dans $\cA\Sigma_0^{-1}$, on pourrait se contenter d'avoir une flèche fonctorielle en sens inverse, du moment qu'elle soit dans $\Sigma_0$, donc puisse s'inverser quand on travaille dans $\cA\Sigma_0^{-1}$. Le candidat evident est $i_X:X\longrightarrow C(X)$ dont on vient de parler. L'ennui, c'est que $i_X$ est seulement dans $\Sigma$, pas dans $\Sigma_0$ (\footnote[42]{S'il n'en était pas ainsi, le type d'homotopie relatif sur $S$ défini par $X/S$ serait localement constant (puisqu'il en est ainsi pour celui défini par $C(X)$), ce qui n'est pas toujours le cas.}). Pourtant, quand $X$ est lui-même de la forme $C(Y)$ (plus généralement, s'il est un Hot-fibré sur $S$), alors le fait que $X\longrightarrow C(X)$ soit dans $\Sigma$ implique déjà qu'il soit dans $\Sigma_0$. Notons donc ceci % [sic] evident -> évident, accent missing
\[
\left\{
\begin{array}{lrl}
i:\id_{\cA}\longrightarrow C\ \text{tel que} & 1^o) & i_X:X\longrightarrow C(X)\ \textit{dans}\ \Sigma\\[1ex]
 & & \text{pour tout }X\in\Ob\,\cA\\[1ex]
 & 2^o) & i_X\ \text{dans }\Sigma_0\text{ si }X\text{ de la}\\[1ex]
 & & \text{forme }C(Y)
\end{array}
\right.
\]

\printedpage{41}
% head: A. Grothendieck : Dérivateurs, Ch. VI, Édition des Universités de Montpellier II et Paris VII

\msp{page 67}
L'idée me vient, après des essais merdouillants \texttt{[sic]}, que $\bar{C}$ pourrait être l'adjoint à droite, et non à gauche \texttt{[c'est la version corrigée (droite et gauche échangées) cf.~p.~65]}, i.e. que les flèches cherchées pourraient aller en sens inverse
\[
\left\{
\begin{aligned}
&\id_{\cA\Sigma^{-1}}&&\xleftarrow{\ p\ }&&\rho\,\bar{C}\quad(=\overline{\rho_0C})\\
&\id_{\cA\Sigma_0^{-1}}&&\xrightarrow{\ i\ }&&\bar{C}\rho\quad(=\overline{C\rho_0})
\end{aligned}
\right.
\]
(j'inverse du coup les notations pour $p$ et $i$), donc on est amené si possible à chercher des flèches
\[
\begin{aligned}
&\id_{\cA}&&\xleftarrow{\ p\ }&&\rho_0C&&=&&C\\
&\id_{\cA}&&\xrightarrow{\ i\ }&&C\rho_0&&=&&C,
\end{aligned}
\]
donc des flèches fonctorielles en $X$
\[
X\xrightarrow{i_X}C(X)\xrightarrow{p_X}X.
\]
On tient bel et bien une $i_X$ fonctorielle - toujours la même! Pour $p_X$, on n'en a pas plus qu'avant, mais se rappelant qu'il suffit de le trouver dans $\cA\Sigma^{-1}$ (et non plus $\cA\Sigma_0^{-1}$!), il suffit d'inverser la flèche $i_X$, qui est en effet un \textit{isomorphisme} dans $\cA\Sigma^{-1}$ (mais non dans $\cA\Sigma_0^{-1}$). On posera donc
\[
p_X=i_X^{-1}:C(X)\longrightarrow X\qquad\quad\text{dans }\cA\Sigma^{-1}.
\]
Il faut prouver deux compatibilités seulement, pour les foncteurs $\bar{C}$ et $\rho$ entre $\cA\Sigma_0^{-1}$ et $\cA\Sigma^{-1}$ (et j'écrirai $C$ au lieu de $\bar{C}$) : les deux composés
\[
\left\{
\begin{array}{l}
\begin{tikzcd}[ampersand replacement=\&]
\rho \arrow[r, "\rho*i"] \arrow[rr, bend left=20, dotted, "\id"] \& \rho\,C\rho \arrow[r, "p*\rho"] \& \rho
\end{tikzcd}
\qquad\quad\text{\footnotesize (relation pour foncteurs à valeurs dans $\cA\Sigma^{-1}$)}\\[2ex]
\begin{tikzcd}[ampersand replacement=\&]
C \arrow[r, "i*C"] \arrow[rr, bend right=20, dotted, "\id"'] \& C\rho\,C \arrow[r, "C*p"] \& C
\end{tikzcd}
\qquad\quad\text{\footnotesize (relation pour foncteurs à valeurs dans $\cA\Sigma_0^{-1}$)}
\end{array}
\right.
\]
doivent être l'identité $\id_\rho$ resp. $\id_C$.

\msp{page 68}
La première relation se vérifie tautologiquement, elle revient à dire que dans $\cA\Sigma^{-1}$, on a un composé identité des
\[
X\xrightarrow{i_X}C(X)\xrightarrow{i_X^{-1}}X\qquad\quad!
\]
La seconde signifie que pour tout $X$, on a un composé identité dans $\cA\Sigma_0^{-1}$
\[
C(X)\xrightarrow{i_{C(X)}}C(C(X))\xrightarrow{C(p_X)}C(X).
\]
Or $p_X:C(X)\longrightarrow X$ avait été défini dans $A\Sigma^{-1}$ comme inverse $i_X^{-1}$ de $i_X$, d'autre part le foncteur $C$ sur $\cA$ définit par hypothèse $\cA\Sigma^{-1}\longrightarrow\cA\Sigma_0^{-1}$, plus précisément $C(\Sigma)\subseteq\Sigma_0$ % [sic] $A\Sigma^{-1}$ (italic A) -> $\mathcal{A}\Sigma^{-1}$, elsewhere script A
\printedpage{42}
% head: A. Grothendieck : Dérivateurs, Ch. VI, Édition des Universités de Montpellier II et Paris VII
ainsi $C(i_X)$ est inversible dans $\cA\Sigma_0^{-1}$, et $C(p_X)=C(i_X)^{-1}$. Donc la compatibilité revient à la relation
\[
C(i_X)^{-1}i_{C(X)}=\id\qquad\quad\textit{dans}\ \ \cA\Sigma_0^{-1},
\]
i.e.
\[
\boxed{\,i_{C(X)}=C(i_X)\,}:C(X)\longrightarrow C^2(X).
\]
En d'autre termes, il faut prouver le : % [sic] d'autre termes -> d'autres termes, plural

\textbf{Lemme.} \textit{Soit $X$ dans $\Cat/S$, considérons les deux flèches}
\[
i_{C(X)},\ C(i_X):C(X)\rightrightarrows C(C(X)).
\]
\textit{Ces flèches sont égales, sinon au sens absolu, du moins dans $(\Cat/S)(W_S^{\mathrm{d}})^{-1}$.}

Ici $f_X:X\longrightarrow S$ est la projection standard et
\[
\begin{aligned}
C(X)\ &=\ \quad\{(x,c)\ \mid\ x\in\Ob\,X,\ c:f_X(x)\ \underbrace{\longrightarrow}_{\text{chemin dans }S}\ t\}\\
C^2(X)\ &=\ \{(x,c,c')\ \mid\ x\in\Ob\,X,\ \underbrace{f_X(x)\xrightarrow{c}t\xrightarrow{c'}t'}_{c,\ c'\text{ chemins dans }S}\},
\end{aligned}
\]
les projections dans $S$ sont données par
\[
\begin{aligned}
(x,c)&&\longmapsto&&t&&=&&\text{but }c\\
(x,c,c')&&\longmapsto&&t'&&=&&\text{but }c'
\end{aligned}
\]

\msp{page 69}
On trouve
\[
i_{C(X)}(x,c)=(x,c,1_t):
\begin{tikzcd}
x \arrow[d] & & \\
f_X(x) \arrow[r, "c"] & t \arrow[r, "1_t", "\text{chemin vide}"'] & t
\end{tikzcd}
\]
\[
C(i_X)(x,c)=(x,1_{f_X(x)},c):
\begin{tikzcd}
x \arrow[d] & & \\
f_X(x) \arrow[r, "1_{f_X(x)}", "\text{chemin vide}"'] & f_X(x) \arrow[r, "c"] & t
\end{tikzcd}
\]
Mais on a le choix entre deux théories de chemins, pour construire $C(X)$ : soit avec les $\ul{Ch}_\infty(S)$, soit avec les $\ul{Ch}(S)$. ‘Ça donne les mêmes résultats’, mais suivant les situations l'un est plus commode que l'autre. Je vais travailler avec les $\ul{Ch}(S)$, auquel cas les $1_S$ sont des unités à gauche et à droite, et on a une loi de composition pour les chemins. Cette loi nous permet de définir
\[
\begin{aligned}
C^2(X)&&\xrightarrow{\ \gamma\ }&&C(X)\\
(x,c,c')&&\longmapsto&&(x,c'\circ c),
\end{aligned}
\]
et les composés de $i_{C(X)}$ et de $C_{i(X)}$ avec cette flèche sont $\id_{C(X)}$ dans les deux cas % [sic] $C_{i(X)}$ -> $C(i_X)$, the arrow considered is C(i_X)
\printedpage{43}
% head: A. Grothendieck : Dérivateurs, Ch. VI, Édition des Universités de Montpellier II et Paris VII
\[
\begin{tikzcd}
C(X) \arrow[r, shift left] \arrow[r, shift right] \arrow[rr, bend right=15, "\id"'] & C^2(X) \arrow[r, "\gamma"] & C(X).
\end{tikzcd}
\]

Donc pour montrer que $i_{C(X)}=C(i_X)$ dans $(\Cat/S)(W_S^{\mathrm{d}})^{-1}$, il suffit de prouver que $\gamma$ est dans $W_S^{\mathrm{d}}$. Mais on sait que le foncteur C transforme $W_S$ en $W_S^{\mathrm{d}}$, donc $C(i_X)$ est dans $W_S^{\mathrm{d}}$, et comme $\gamma\circ C(i_X)=\id$, $\gamma$ est dans $W_S^{\mathrm{d}}$, OK! % [sic] foncteur C (upright roman C) -> foncteur $C$, math italic elsewhere

\msp{page 70}
J'ai envie de grouper ce qui a été prouvé dans un énoncé récapitulatif.

\textbf{Théorème-Scholie.} On considère le diagramme de catégories de la page 8 (\footnote[43]{On suppose sur $W$ les conditions a) b) c) de p.~45, de plus L5 (p.~43). Cf. commentaires p.~81.}).

\circled{I} Quand on munit toutes les 13 catégories du diagramme du localiseur $W_S$ ($W$-équivalence grossière, i.e. $W$-équivalence dans Cat, sans plus), \textit{tous} les foncteurs du diagramme induisent des équivalences sur les catégories localisées (\footnote[44]{De plus, si $\mathcal{F}$ est une sous-catégorie pleine de $\Cat/S$ qui contient une des deux sous-catégories $\ul{Fib}_0S$, $\ul{Cofib}_0S$ (ce qui inclut le cas des 8 catégories des diagrammes, autres que $\Cat/S$ elle-même), l'inclusion $(\mathcal{F},W_S|\mathcal{F})$ dans $(\Cat/S,W_S)$ induit une équivalence pour les localisés.}). Désignons par
\[
\HOT^{\mathrm{lc}}_W(S)
\]
‘la’ catégorie localisée commune.

\circled{II} Munissons les quatre catégories bordantes gauches
\[
(*)\qquad\qquad
\begin{tikzcd}
\Cat/S & \ul{Liss}\,S \arrow[l, hook'] & \ul{Fib}\,S \arrow[l, hook'] & \ul{Fibsc}\,S \arrow[l, hook']
\end{tikzcd}
\]
du localiseur $W_S^{\mathrm{d}}$ de la $W$-équivalence fine \textit{droite} sur $S$ ( $\begin{tikzcd}[row sep=small, column sep=small]
X \arrow[rr, "f"] \arrow[dr] & & Y \arrow[dl]\\
 & S &
\end{tikzcd}$ est dans $W_S^{\mathrm{d}}$ si et seulement si les $s\backslash X\longrightarrow s\backslash Y$ le sont \texttt{[plutôt, sss ils sont dans $W$]}, pour tout $s\in\Ob\,S$). Les foncteurs du diagramme, i.e. ceux de $(*)$, entre ces catégories, donc aussi leurs composés, induisent des équivalences entre les localisées (\footnote[45]{Plus généralement, ça marche pour toute sous-catégorie pleine de $\Cat/S$ qui contient $\ul{Fib}\,S$ (et dualement, avec $W_S^{\mathrm{g}}$ au lieu de $W_S^{\mathrm{d}}$, si $\mathcal{F}$ contient $\ul{Cofib}\,S$).}). Désignons par
\[
\HOT_W(S)
\]
‘la’ catégorie localisée commune.

Enoncé dual pour
\[
(*')\qquad\qquad
\begin{tikzcd}
\Cat/S & \ul{Prop}\,S \arrow[l, hook'] & \ul{Cofib}\,S \arrow[l, hook'] & \ul{Cofibsc}\,S, \arrow[l, hook']
\end{tikzcd}
\]

% foot (each page): Version 9 novembre 2001

\end{document}
